\chapter{Data Validation as an Executable Contract}
\label{ch:data-validation}

\begin{learningobjectives}
After completing this chapter, you will be able to:
\begin{itemize}
  \item distinguish validation from cleaning and profiling;
  \item classify source, key, relationship, range, category, and temporal checks;
  \item write duplicate, orphan, range, and category validation queries;
  \item interpret error, warning, and informational severity;
  \item explain why zero failing rows is meaningful only with a stated rule; and
  \item use the project's persisted quality results to decide whether analysis
    may continue.
\end{itemize}
\end{learningobjectives}

\section{Validation is not decoration}

\term{Data validation} compares observed data with an explicit contract.
\term{Data cleaning} transforms values or records according to a documented
rule. \term{Profiling} describes the data without necessarily declaring it
valid or invalid.

\begin{plainlanguage}
Profiling says ``1,111 values are missing.'' Validation says whether that
condition violates a rule. Cleaning says what transformation, if any, should
follow. The three steps answer different questions.
\end{plainlanguage}

The project persists every SQL check in \texttt{quality.run\_results} with:
\begin{itemize}
  \item run timestamp;
  \item check name;
  \item severity;
  \item number of failing rows;
  \item status; and
  \item human-readable detail.
\end{itemize}

\section{Severity}

\begin{table}[htbp]
\centering
\caption{Implemented quality severity contract}
\label{tab:severity}
\begin{tabularx}{\textwidth}{L{0.18\textwidth}L{0.29\textwidth}Y}
\toprule
Severity & Meaning in this project & Example \\
\midrule
Error & A violated condition blocks a trusted pipeline run & Duplicate attempt keys or invalid scores \\
Warning & An unexpected condition deserves review but may be explainable & Unexpected module or presentation code \\
Informational & A known condition is profiled and kept visible & Missing deprivation band or pre-course activity \\
\bottomrule
\end{tabularx}
\end{table}

An informational profile is not a silent pass. It records a real feature of the
data whose analytical consequences must be considered.

\section{Source-count checks}

The project records fixed row-count expectations for especially important
sources:

\begin{lstlisting}[language=SQL,caption={Simplified form of an implemented
source-count check.}]
SELECT
    abs((SELECT count(*) FROM raw.courses) - 22)
        AS failing_rows,
    CASE
        WHEN (SELECT count(*) FROM raw.courses) = 22
        THEN 'pass'
        ELSE 'fail'
    END AS status;
\end{lstlisting}

The analytical question is: ``Did the imported source match the verified
archive?'' The input grain is one course row per module-presentation. The
output grain is one check result. A nonzero difference is an error because the
database no longer reconciles with the expected archive.

\section{Natural-key checks}

To test the attempt key:

\begin{lstlisting}[language=SQL]
SELECT count(*) AS duplicate_keys
FROM (
    SELECT
        code_module,
        code_presentation,
        id_student
    FROM staging.student_info
    GROUP BY 1, 2, 3
    HAVING count(*) > 1
) AS duplicated;
\end{lstlisting}

The inner query returns one row per duplicated key, not one row per duplicated
source record. The outer count therefore reports duplicated key groups.

\begin{validationbox}
Always ask what \texttt{failing\_rows} counts. It might count:
\begin{itemize}
  \item bad source rows;
  \item duplicated key groups;
  \item orphan child rows; or
  \item the absolute difference from an expected count.
\end{itemize}
The detail column should make the unit explicit.
\end{validationbox}

\section{Relationship checks with anti-joins}

An \term{orphan} is a child row whose required parent is missing. The project
uses a left anti-join pattern:

\begin{lstlisting}[language=SQL]
SELECT count(*) AS orphan_submissions
FROM staging.assessment_submissions AS submission
LEFT JOIN staging.assessments AS assessment
  USING (id_assessment)
WHERE assessment.id_assessment IS NULL;
\end{lstlisting}

\subsection{Reasoning through the query}

\begin{enumerate}
  \item The left input is one recorded student-assessment submission.
  \item The right input is one assessment.
  \item The join key is \texttt{id\_assessment}.
  \item A matching assessment supplies a nonnull right-side ID.
  \item The \texttt{WHERE} clause retains only unmatched submissions.
  \item The final output is one count.
\end{enumerate}

The expected count is zero. A nonzero result means a submission references an
assessment absent from the project source model.

\section{Range and category checks}

\begin{lstlisting}[language=SQL,caption={Implemented validation patterns for
scores and outcomes.}]
SELECT count(*) AS invalid_scores
FROM staging.assessment_submissions
WHERE score NOT BETWEEN 0 AND 100;

SELECT count(*) AS invalid_results
FROM staging.student_info
WHERE final_result NOT IN (
    'Distinction', 'Pass', 'Fail', 'Withdrawn'
);
\end{lstlisting}

\texttt{BETWEEN 0 AND 100} includes both boundaries. \texttt{NOT IN} is useful
for a documented finite set, but null values need separate reasoning because
SQL's three-valued logic does not classify null as ``not in'' the set.

\section{Date and temporal checks}

The raw dates are course-relative integers. The project checks that
unregistration does not precede registration:

\begin{lstlisting}[language=SQL]
SELECT count(*) AS invalid_registration_order
FROM staging.student_registration
WHERE unregistration_day < registration_day;
\end{lstlisting}

The later snapshot test checks that the latest activity used in a feature is no
later than the snapshot:

\begin{lstlisting}[language=SQL]
SELECT count(*) AS future_information
FROM features.model_snapshots
WHERE latest_feature_day > snapshot_day;
\end{lstlisting}

The second query previews the forthcoming leakage-prevention chapter. It is
shown here because it is part of the real 24-check quality result, not because
time-aware modeling is fully taught in Edition 0.1.

\section{Known profiles are evidence}

The saved project run reports:

\begin{table}[htbp]
\centering
\caption{Selected persisted quality profiles}
\label{tab:quality-profiles}
\begin{tabularx}{\textwidth}{L{0.38\textwidth}L{0.18\textwidth}Y}
\toprule
Check & Failing/profiled rows & Interpretation \\
\midrule
Error-level checks & 0 across 18 checks & The recorded error contracts passed \\
Activity before course start & 688,988 & Legitimate negative relative days; excluded from model weeks \\
Missing deprivation band & 1,111 & Missing context remains visible \\
Missing assessment due dates & 11 & Often exams; excluded from next-assessment targets \\
Unexpected module codes & 0 & Codes stayed within AAA--GGG \\
Unexpected presentation codes & 0 & Codes matched documented year/term format \\
\bottomrule
\end{tabularx}
\end{table}

\verified{The saved SQL quality table contains 24 results: 18 error-level
checks, three warnings, and three informational profiles. No error-level check
failed in that run.}

\section{Run the quality gate}

After raw import and the ordered SQL pipeline:

\begin{lstlisting}[language=bash]
make sql
make validate
\end{lstlisting}

\projectfile{scripts/run_quality.sh} reruns source and relational checks. If the
snapshot table exists, it also runs the four point-in-time checks. It then
counts failed error-level results and exits with a failure status when the count
is nonzero.

\begin{validationbox}
The gate passes only when:
\begin{enumerate}
  \item the shell command exits successfully;
  \item every error-level row has status \texttt{pass};
  \item warnings and informational profiles have been read;
  \item the run timestamp is current for the database under study; and
  \item no expected check silently disappeared.
\end{enumerate}
\end{validationbox}

\section{How to investigate a failed check}

Use a consistent procedure:
\begin{enumerate}
  \item Identify the check name and what its failing count represents.
  \item Open the implementing SQL and isolate the detail query.
  \item Replace the outer count with selected identifying columns.
  \item Inspect a small bounded sample without editing source data.
  \item Compare with official documentation and source audit.
  \item Decide whether the issue is import, source, representation, or rule.
  \item Document the decision before any cleaning change.
\end{enumerate}

\begin{editionnote}
This course does not change project quality rules. If a rule appears incomplete
or incorrect, record the evidence as a suggested enhancement. Do not silently
weaken the check to make the pipeline pass.
\end{editionnote}

\section{Knowledge check}

\begin{enumerate}
  \item Distinguish profiling, validation, and cleaning.
  \item What does an anti-join find?
  \item Why can \texttt{NOT IN} require a separate null check?
  \item What does zero mean in the duplicate-attempt query?
  \item Why must warnings be read even when the gate exits successfully?
  \item What is the first step after a failed quality check?
\end{enumerate}

\begin{exercisebox}{Exercise 7.1: Write a validation specification}
For each rule below, write its severity, input grain, failing-row meaning,
expected result, and investigation query:
\begin{enumerate}
  \item assessment weights are between 0 and 100;
  \item every VLE interaction site exists in the same module-presentation;
  \item no attempt key is duplicated;
  \item deprivation band may be missing but must be profiled; and
  \item unregistration cannot precede registration.
\end{enumerate}
\solutionref{sol:validation-spec}
\end{exercisebox}

\begin{exercisebox}{Exercise 7.2: Debug an orphan count}
Suppose the orphan-submission check returns 12. Explain how to turn the count
query into a diagnostic query, which identifiers to display, which source files
to inspect, and why deleting the 12 rows is not an acceptable first response.
\solutionref{sol:orphan-debug}
\end{exercisebox}

\begin{commonmistakes}
\begin{itemize}
  \item Treating every nonzero profile as an error.
  \item Treating a zero count as meaningful without understanding its unit.
  \item Changing source data before deciding whether the rule or import is
    wrong.
  \item Ignoring null behavior in range or membership conditions.
  \item Reporting an old persisted quality run as evidence for a new database.
\end{itemize}
\end{commonmistakes}

\transferheading
Build validation around contracts that name a grain, rule, severity, failing
unit, and expected response. Persist results with timestamps. This pattern
works for event logs, financial ledgers, clinical records, and any other
multi-table analytical system.

\begin{summarybox}
Validation turns assumptions into executable, persisted evidence. The project
checks source counts, natural keys, relationships, ranges, categories, dates,
and future-information constraints. Error, warning, and informational results
have different meanings. A successful gate is useful only when the checks,
failing units, timestamps, and known profiles are understood.
\end{summarybox}
